\lyseis{Β3.8}{Γωνίες σε παράλληλες ευθείες με τέμνουσα}
\sol{1} α) εντός--εκτός και επί τα αυτά. β) εντός εναλλάξ. γ) εκτός εναλλάξ. δ) εντός και επί τα αυτά. ε) εντός και επί τα αυτά. στ) εντός--εκτός και επί τα αυτά.
\sol{2} Οι οξείες: $\alpha=\gamma=\varepsilon=\eta=70^\circ$ (κατακορυφήν, εντός--εκτός, εκτός εναλλάξ). Οι αμβλείες: $\beta=\delta=\zeta=\theta=180^\circ-70^\circ=110^\circ$.
\sol{3} $\beta=118^\circ$ (κατακορυφήν με τη $\delta$). $\zeta=118^\circ$ (εντός εναλλάξ με τη $\delta$). $\theta=118^\circ$ (εντός--εκτός και επί τα αυτά με τη $\delta$). $\varepsilon=180^\circ-118^\circ=62^\circ$ (εντός και επί τα αυτά με τη $\delta$). $\alpha=62^\circ$ (παραπληρωματική της $\beta$).
\sol{4} Οι εντός εναλλάξ είναι ίσες: $3x-10^\circ=x+50^\circ$, άρα $2x=60^\circ$ και $x=30^\circ$. Κάθε γωνία είναι $80^\circ$.
\sol{5} Παραπληρωματικές: $x+4x=180^\circ$, άρα $5x=180^\circ$ και $x=36^\circ$. Οι γωνίες είναι $36^\circ$ και $144^\circ$.
\sol{6} α) Λάθος, ισχύει μόνο για παράλληλες. β) Σωστό. γ) Λάθος, είναι παραπληρωματικές. δ) Σωστό. ε) Σωστό (αντίστροφο).
\sol{7} Ελεύθερη κατασκευή. Πρέπει να βρεις τέσσερις γωνίες $55^\circ$ και τέσσερις $125^\circ$ και να επιβεβαιώσεις με τον μοιρογνωμόνιο ότι $55^\circ+125^\circ=180^\circ$.
\sol{8} Τέσσερις γωνίες είναι $62^\circ$ και τέσσερις $118^\circ$.
\sol{9} Η εντός εναλλάξ είναι ίση: $48^\circ$. Η εντός και επί τα αυτά είναι παραπληρωματική: $180^\circ-48^\circ=132^\circ$.
\sol{10} Οι γωνίες είναι εντός και επί τα αυτά, άρα παραπληρωματικές: $180^\circ-125^\circ=55^\circ$.
\sol{11} Όχι. Αν τα δοκάρια ήταν παράλληλα, οι εντός εναλλάξ γωνίες θα ήταν ίσες, αλλά $74^\circ\neq76^\circ$.
\sol{12} Ναι. Οι εντός και επί τα αυτά γωνίες έχουν άθροισμα $98^\circ+82^\circ=180^\circ$, άρα οι πλευρές είναι παράλληλες (αντίστροφο).
\sol{13} Έστω ότι $\gamma=\varepsilon$ (εντός εναλλάξ). Η $\varepsilon$ και η $\zeta$ βρίσκονται πάνω στην ευθεία $\varepsilon_2$ και σχηματίζουν ευθεία γωνία, άρα $\varepsilon+\zeta=180^\circ$. Επειδή $\gamma=\varepsilon$, είναι $\gamma+\zeta=180^\circ$, δηλαδή οι εντός και επί τα αυτά γωνίες είναι παραπληρωματικές.
\sol{14} $\widehat{\mathrm{ΒΕΓ}}=40^\circ+30^\circ=70^\circ$. Η παράλληλη από το $\mathrm{Ε}$ χωρίζει τη γωνία σε δύο μέρη. Το πρώτο είναι ίσο με $40^\circ$ (εντός εναλλάξ με τη γωνία στο $\mathrm{Β}$) και το δεύτερο ίσο με $30^\circ$ (εντός εναλλάξ με τη γωνία στο $\mathrm{Γ}$).
\sol{15} Οι γωνίες του τετραπλεύρου είναι $70^\circ,110^\circ,70^\circ,110^\circ$. Οι διαδοχικές είναι εντός και επί τα αυτά, άρα παραπληρωματικές, και οι απέναντι ίσες.
\sol{16} Αν η οξεία είναι $x$, τότε $x+3x=180^\circ$, άρα $x=45^\circ$. Οι γωνίες είναι $45^\circ$ και $135^\circ$.
\sol{17} Ελεύθερη απάντηση. Οι εντός εναλλάξ γωνίες πρέπει να βγουν ίσες (με μικρή απόκλιση λόγω μέτρησης).
\sol{18} Ελεύθερη κατασκευή. Σε κάθε σημείο τομής θα έχεις δύο ζεύγη ίσων οξειών και δύο ζεύγη ίσων αμβλειών, και μια οξεία με μια αμβλεία έχουν άθροισμα $180^\circ$.
